\section{Wave 13 Nekrasov adversarial attack-heal on the DT side of
Conjecture 1: roots of $Z^X_{L,h_M}$, the eight Gritsenko--Cl\'ery
forms, and twined K3 elliptic genera}
\label{wn:wave13-a5-dt-zeta-nekrasov}

\providecommand{\ClaimStatusTheorem}{\textsc{[theorem]}}
\providecommand{\ClaimStatusConjectured}{\textsc{[conjectural]}}
\providecommand{\ClaimStatusDefinition}{\textsc{[definition]}}
\providecommand{\ClaimStatusDerivation}{\textsc{[first-principles]}}
\providecommand{\ClaimStatusHeuristic}{\textsc{[heuristic]}}
\providecommand{\kBKM}{\kappa_{\mathrm{BKM}}}
\providecommand{\kch}{\kappa_{\mathrm{ch}}}
\providecommand{\kcat}{\kappa_{\mathrm{cat}}}
\providecommand{\kfib}{\kappa_{\mathrm{fiber}}}
\providecommand{\Mtf}{M_{24}}
\providecommand{\ZZ}{\mathbb{Z}}
\providecommand{\CC}{\mathbb{C}}
\providecommand{\HH}{\mathbb{H}}
\providecommand{\QQ}{\mathbb{Q}}
\providecommand{\Sp}{\mathrm{Sp}}
\providecommand{\SO}{\mathrm{SO}}
\providecommand{\Aut}{\mathrm{Aut}}
\providecommand{\Phiten}{\Phi_{10}}
\providecommand{\AAtwo}{\mathcal{A}_2}
\providecommand{\fg}{\mathfrak{g}}

Five pairs of contradictions $(A_i\mid H_i)$ on the DT-side of
Conjecture~1 of Lorgat 2020 (source paper lines 96--121; slides
lines 911--936). Residue, Humbert divisor, Oberdieck--Pixton
denominator, twined elliptic genus, Borcherds product. Nothing else.
For each attack, the ghost theorem is extracted, the linguistic slip
named, and the correct relationship stated on the domain where the
statement is actually true --- the Siegel upper half plane $\HH_2$
covering the Siegel modular threefold $\AAtwo = \Sp_4(\ZZ)\backslash
\HH_2$, where $\Delta_5$ and its seven Gritsenko--Cl\'ery congeners
live as paramodular cusp forms and where the DT partition
$Z^X_{L,h_M}$ pulls back via Igusa's isomorphism
$\Sp_4(\ZZ)/\{\pm I_5\}\cong\SO(\Lambda^{3,2})_+$ (Lemma~1 of Lorgat
2020). Scope of the Oberdieck--Pandharipande theorem is enforced at
every step: primitive $\beta\in H_2(K3,\ZZ)$, $N=1$ is theorem; the
other seven rows are conjectural. Nekrasov voice throughout; no
$\Omega$-background shortcut (K3 is not toric), three independent
verification paths at $N=1$, separate treatment of $N\in\{2,3,4,6\}$
and the remaining $N=8$ and the double-index $\Gamma_1(2), \Gamma_2(2),
\Gamma_1(3), \Gamma_1(4)$ rows.

%--------------------------------------------------------------------------
\subsection*{A1. ``DT zeta roots'': residue of a rational function,
or pole locus of a Siegel modular form?}

\textsc{Attack.} The phrase ``DT zeta roots'' in Conjecture~1 of
Lorgat 2020 (paper line 116: ``all eight diagonal divisor modular
forms of Gritsenko--Cl\'ery arise, up to constant $C$, as
reciprocal-square roots of $Z^X_{L,h_M}$'') invites four distinct
readings:
\begin{enumerate}
\item \emph{Functional roots} of $Z^X_{L,h_M}(q,t,p)$ as a meromorphic
    function of three variables: points $(q_0,t_0,p_0)$ where
    $Z^X_{L,h_M}=0$.
\item \emph{Denominator poles} of $Z^X_{L,h_M}$ as a rational
    function of the DT generating variables: zero-divisor of the
    reciprocal $1/Z^X_{L,h_M}$.
\item \emph{Reciprocal-square roots} in the literal sense of the
    Conjecture 1 text: the eight Gritsenko--Cl\'ery forms $F_i$
    realise $Z^X_{L,h_M} = C\cdot F_i^{-2}$, so that $F_i =
    \pm\sqrt{C}/\sqrt{Z^X_{L,h_M}}$ is an eighth-of-a-square-root
    extraction.
\item \emph{Humbert-divisor locus}: since $\Delta_5$ vanishes along
    Humbert divisors $H_n\subset\AAtwo$ of discriminant $n\in
    1+4\ZZ_{\ge 0}$ (Gritsenko--Hulek--Sankaran 2007, for $\Delta_5$
    only to order exactly one), the ``root locus'' of $F_i$ is the
    paramodular pull-back of a union of Humbert divisors.
\end{enumerate}

\noindent
\textbf{Ghost theorem.} Reading (3) is literally the paper and
Oberdieck--Pixton 2018 (arXiv:1706.10100 Theorem~1): at $N=1$,
\[
Z^{S\times E}(q,t,p)
  \;=\;-\,\frac{C}{\Phiten(\tau,z,\sigma)}
  \;=\;-\,\frac{C}{\Delta_5(\tau,z,\sigma)^2},
\]
where $(q,t,p)=(e^{2\pi i\tau},e^{2\pi iz},e^{2\pi i\sigma})$ and the
minus sign is the KKV/MNOP/MPT sign from the reduced-virtual-class
$(-1)^{\mathrm{vdim}}$ convention. So reading (3) makes Conjecture 1
precisely: \emph{$F_i^2\cdot Z^X = -C$, i.e.\ $F_i$ is the
paramodular denominator of the reduced DT partition.} Reading (4),
\emph{the vanishing locus of $F_i$ in $\AAtwo$ is a Humbert divisor
decomposition}, is the Borcherds-product-side realisation.

\noindent
\textbf{Precise error.} Reading ``DT zeta roots'' as reading (1) --- a
set of zeroes of $Z^X$ --- is wrong: $Z^X$ is everywhere nonzero
generically on $\HH_2$, with isolated \emph{poles} along Humbert
divisors coming from $\Delta_5$'s zero locus (singular abelian
surfaces: product of two elliptic curves $E_1\times E_2$ with shared
level structure). The ``root'' in the phrase ``DT zeta root'' is a
\emph{second root} (reciprocal square root), not a \emph{zero of $Z^X$}.
The two readings coincide only on the pole locus, which is reading
(4).

\noindent
\textbf{Correct relationship.} The pole divisor of
$1/Z^{S\times E}=-\Delta_5^2/C$ on the Igusa threefold
$\AAtwo^*=\AAtwo\cup$(cusps) decomposes as a union of Humbert
divisors $H_n$. For $\Delta_5$ (Gritsenko--Nikulin 1998;
Gritsenko--Hulek--Sankaran 2007 Theorem~1.1),
\[
\mathrm{div}(\Delta_5)\;=\;H_1\qquad\text{(as divisor on }\AAtwo^*),
\]
with vanishing order exactly $1$ along $H_1$ --- the ``simplest
diagonal divisor'' of Question~1. In particular,
$\mathrm{div}(1/Z^{S\times E})=2H_1$, so the ``DT roots'' of
Conjecture~1, read as the divisor where the DT partition diverges,
are twice the simplest Humbert divisor. Each of the other seven
Gritsenko--Cl\'ery forms $F_i$ has divisor $H_{n_i}$ for a specific
$n_i\in 1+4\ZZ_{\ge 0}$ on its paramodular cover of $\AAtwo^*$, and
Conjecture~1 reads: for each $N\in$ scope-set, there exists a
polarised K3 datum $(S,L,g_N,h_M)$ with
$Z^X_{L,h_M}=-C/F_i^2$ and $\mathrm{div}(F_i)=H_{n_i}$ on the
paramodular Siegel threefold.

\textsc{Heal.} In working-notes chapter
\texttt{chapters/examples/k3e\_bkm\_chapter.tex} and the Lorgat 2020
cite in bibliography, replace ``DT zeta roots'' (paper-line 116
language) with the precise triple
\emph{(reciprocal-square-root of $Z^X$; denominator divisor
$H_{n_i}\subset\AAtwo^*$; vanishing order one)}. Trace the
three-variable $(q,t,p)\leftrightarrow(\tau,z,\sigma)$ dictionary
(Igusa 1964 coordinates on $\HH_2$) and the primitive-class
restriction (OP 2017 p.~2; cf.\
\texttt{wave13\_f4\_oberdieck\_primitive\_scope.tex} Cycle~2).

%--------------------------------------------------------------------------
\subsection*{A2. The eight Gritsenko--Cl\'ery forms: a complete list
with weight, character, paramodular group, and Humbert discriminant}

\textsc{Attack.} The slides line 50--54 give the eight diagonal-divisor
forms as
\[
M_5(\Gamma_1,\nu_2)\;=\;\Delta_5,\qquad
M_2(\Gamma_2,\nu_4),\qquad
M_3(\Gamma_1(2),\nu_2),
\]
\[
M_1(\Gamma_3,\nu_6),\qquad M_2(\Gamma_1(3),\nu_2),
\]
\[
M_{1/2}(\Gamma_4,\nu_8),\qquad
M_{3/2}(\Gamma_1(4),\nu_4),\qquad
M_1(\Gamma_2(2),\nu_4),
\]
where $M_k(\Gamma,\nu)$ denotes a weight-$k$ modular form with
character $\nu$ on the paramodular-type group $\Gamma$. This list,
as stated, is opaque: the paramodular-group data, the multiplier
orders, and the connection to $N$-torsion CY$_3$ data are all
implicit. Worse: claiming ``exactly eight'' (Theorem 1, paper line
56) without pinning which of the $(\Gamma_t,\nu_j)$ subgroups the
scope is over admits multiple inconsistent readings. And the
physics labelling in the main paper (paper line 112: ``$N\le 8$ by
Nikulin'') collides with the apparent symplectic-auto order
$N\in\{1,2,3,4\}$ implicit in the first five entries, plus the
auxiliary half-level $\Gamma_t(N)$ Hecke data.

\noindent
\textbf{Ghost theorem.} Gritsenko--Cl\'ery
(\emph{arXiv:0812.3962}, 2008) Theorem~1: for the paramodular groups
$\Gamma_t(N)<\Gamma_t$ ($t,N\in\ZZ_{>0}$, $t\le 4$, $N|2t$), there
are \emph{exactly eight} Siegel modular forms with character that
vanish to order exactly $1$ along the diagonal of $\HH_2$; they are
listed above, and each has a \emph{Borcherds-Gritsenko lift} from a
weak Jacobi form of weight $0$ and index $t$ (the four even-index
rows), respectively index $t/2$ (the four half-index rows).

\noindent
\textbf{Precise error prevented.} Confusing the ``$8$'' of
Gritsenko--Cl\'ery (over all paramodular subgroups of $\Sp_4(\QQ)$
with the diagonal-divisor-of-order-one property) with the ``$8$'' of
Nikulin (maximal order of a symplectic automorphism of a K3 surface,
$\mathrm{ord}(g)\le 8$). They are the same number but different
theorems: Gritsenko--Cl\'ery is an exhaustion over paramodular
groups; Nikulin is a topological constraint on K3 automorphisms. The
identification of an 8-form with an $N$-torsion CY$_3$ datum is
\emph{Conjecture 1} of Lorgat 2020, the content of the whole
discussion.

\noindent
\textbf{Correct relationship (eight-row explicit table).} In the
basis $(q,\zeta,q')=(e^{2\pi i\tau},e^{2\pi iz},e^{2\pi i\sigma})$
with $(\tau,z,\sigma)\in\HH_2$, the eight Gritsenko--Cl\'ery forms
stratify as follows.

\begin{center}
\renewcommand{\arraystretch}{1.3}
\begin{tabular}{llllll}
\toprule
$i$ & Form & Weight $k_i$ & Group $\Gamma_{t_i}(N_i)$ &
Character ord.\ $\#\nu_i$ & Humbert disc.\ $n_i$\\
\midrule
1 & $\Delta_5$                 & $5$   & $\Gamma_1=\Sp_4(\ZZ)$    & $2$ & $1$ \\
2 & $\Delta_2^{(2)}$          & $2$   & $\Gamma_2$               & $4$ & $4$ \\
3 & $\Delta_3^{(1,2)}$         & $3$   & $\Gamma_1(2)$            & $2$ & $4$ \\
4 & $\Delta_1^{(3)}$          & $1$   & $\Gamma_3$               & $6$ & $9$ \\
5 & $\Delta_2^{(1,3)}$        & $2$   & $\Gamma_1(3)$            & $2$ & $9$ \\
6 & $\Delta_{1/2}^{(4)}$      & $1/2$ & $\Gamma_4$               & $8$ & $16$ \\
7 & $\Delta_{3/2}^{(1,4)}$    & $3/2$ & $\Gamma_1(4)$            & $4$ & $16$ \\
8 & $\Delta_1^{(2,2)}$        & $1$   & $\Gamma_2(2)$            & $4$ & $4$ \\
\bottomrule
\end{tabular}
\end{center}

The Humbert discriminant $n_i$ is the smallest $n$ such that $H_n$
appears in the divisor of $F_i$ on the paramodular Siegel threefold.
The character orders $\#\nu_i\in\{2,4,6,8\}$ factor through
$\mathrm{GL}_2(\ZZ)/\ker(\nu_i)$ and the Maass explicit multiplier
system for $\Gamma_t(N)$ (a direct extension of the $\nu_{\Delta_5}$
formula from paper line 144--164).

\emph{Three independent verifications at $N=1$
(row $i=1$: $\Delta_5$).}

(i)~\emph{Infinite-product Borcherds lift.} Theorem~4 of
Lorgat 2020:
\[
\tfrac{1}{64}\Delta_5(\tau,z,\sigma)\;=\;e^{\pi i(\tau+z+\sigma)}
  \prod_{\substack{n,\ell,m\in\ZZ\\(n,\ell,m)>0}}
  (1-q^n\zeta^\ell q'^m)^{f(nm,\ell)},
\]
with $f(nm,\ell)$ the Fourier coefficients of $\phi_{0,1}$ (paper
line 716--724). At the cusp $T_0=\bigl(\begin{smallmatrix}
1&1/2\\1/2&1\end{smallmatrix}\bigr)$ the product is
$1+O(e^{-\varepsilon})$; weight $5$ read off the transformation
$\Delta_5(Z+B)=(-1)^{b_1+b_2+b_3}\Delta_5(Z)$ (paper line 156).

(ii)~\emph{Gritsenko additive lift.} $\Delta_5$ is the additive
Gritsenko lift of $\phi_{5,1/2}=\eta^9\vartheta_{11}$, the unique
Jacobi cusp form of weight $5$ and index $1/2$ with the
$\nu_{\Delta_5}$-character (paper line 204--226).

(iii)~\emph{Igusa ring-generator relation.} $\Delta_5^2=\Delta_{10}
\in\CC[E_4,E_6,\Delta_{10},\Delta_{12}]=\SM(\Sp_4(\ZZ))$, with
$\Delta_{10}=\Phiten$ the Igusa cusp form; Igusa 1964 identifies
$\Delta_5$ as the unique (up to scalar) square root of the weight-10
generator.

All three identify the same modular form.

\textsc{Heal.} Insert the above eight-row table in
\texttt{chapters/examples/k3e\_bkm\_chapter.tex}
\S``Gritsenko--Cl\'ery catalogue''. Attach to each row the
three-path verification of row $i=1$ (the $N=1$ case, which is the
only theorem; rows 2--8 are conjectures pending the conjectural
twisted-DT theorems below). Flag all rows $i\ge 2$
\ClaimStatusConjectured, pointing to the conjectural
Oberdieck--Pandharipande 2014 Conjecture~B extension
(\texttt{wave13\_f4\_oberdieck\_primitive\_scope.tex} Cycle~3).

%--------------------------------------------------------------------------
\subsection*{A3. Mathieu correspondence: twined elliptic genera
$Z_g(\tau,z)$ for $g\in\Mtf$ and their role in the eight-form map}

\textsc{Attack.} Conjecture~1 says (paper line 117--118):
\emph{``these Siegel paramodular forms all arise as denominator
functions of generalized Borcherds-Kac-Moody superalgebras, with
root multiplicities specified by $g_N$-$h_M$-twisted twined elliptic
genera of K3 surfaces.''} Two layers of ambiguity remain:
\begin{itemize}
\item Mathieu moonshine (Eguchi--Ooguri--Tachikawa 2010,
    Gaberdiel--Hohenegger--Volpato 2012) maps symplectic K3
    automorphisms $g$ of finite order to conjugacy classes
    $[g]\in\Mtf$ via the \emph{twining}
    $Z_g(\tau,z)=\mathrm{Tr}_{H(K3)}\,g\,y^{F_L}q^{L_0-c/24}$,
    with $Z_g$ a weak Jacobi form of weight $0$ and index $1$ for the
    congruence subgroup $\Gamma_0(\mathrm{ord}(g))$ times a
    prescribed multiplier (GHV 2012 Theorem~3.1). But
    \emph{which} conjugacy class matches \emph{which}
    Gritsenko--Cl\'ery row is not fixed by the conjecture.
\item The double-twisting $(g_N,h_M)$ in the paper invokes a
    \emph{pair} of symplectic automorphisms; the Mathieu literature
    handles primarily single-twining $Z_g$ and, for orbifolds,
    single-shift $g$-orbifolded $Z_{g|h}$. The $(g_N, h_M)$-twined
    case is post-GHV 2012 and is the content of
    Persson--Volpato 2014 (\texttt{arXiv:1312.6500}, umbral
    subconnection) and
    Bryan--Oberdieck--Pandharipande--Yin 2018
    (\texttt{arXiv:1802.09180}, CHL/CY$_3$ twining tables).
\end{itemize}

\noindent
\textbf{Ghost theorem.} Bryan--Oberdieck--Pandharipande--Yin 2018
Theorem~7 (``Twisted K3$\times$E correspondence''): for any
$N\in\{1,2,3,4,6\}$ and primitive $\beta\in H_2(K3,\ZZ)^{g_N}$, the
twisted Noether--Lefschetz DT partition
$Z^{(S\times E)/\langle g_N\rangle}_{\mathrm{prim}}(\tau,z,\sigma)$ equals
$-1/F_N(\tau,z,\sigma)^2$, where $F_N$ is a paramodular cusp form with
Borcherds lift from the twined K3 elliptic genus
$Z_{g_N}(\tau,z)\in J_{0,1}(\Gamma_0(N),\psi_N)$.
\emph{This theorem is established for $N\in\{2,3,4,6\}$
conditionally on Conjecture~B of Oberdieck--Pandharipande 2014 at
imprimitive classes.}

\noindent
\textbf{Precise error.} Two conjectures quietly conflated: (a) that
the map $N\mapsto F_N$ is a bijection onto five of the eight
Gritsenko--Cl\'ery rows; (b) that the $(g_N,h_M)$-double-twining
(with $h_M$ an auxiliary order-$M$ symplectic auto commuting with
$g_N$) generates the remaining three rows. Neither is theorem.

\noindent
\textbf{Correct relationship (five-plus-three mapping).} The
Mathieu-$\Mtf$-conjugacy-to-row map is as follows (rows labelled
as in A2 table).

\begin{center}
\renewcommand{\arraystretch}{1.3}
\begin{tabular}{llll}
\toprule
Row $i$ & $(N,M)$ & $\Mtf$-class $[g_N]$ (Frame shape) & Conjecture-1
status\\
\midrule
$1$ & $(1,-)$ & $1A$ $(1^{24})$                 & \ClaimStatusTheorem\\
$2$ & $(2,-)$ & $2A$ $(1^8 2^8)$                 & \ClaimStatusConjectured\\
$3$ & $(1,2)$ & $2A$ (shifted) $(1^8 2^8)$      & \ClaimStatusConjectured\\
$4$ & $(3,-)$ & $3A$ $(1^6 3^6)$                 & \ClaimStatusConjectured\\
$5$ & $(1,3)$ & $3B$ $(3^8)$                    & \ClaimStatusConjectured\\
$6$ & $(4,-)$ & $4B$ $(1^4 2^2 4^4)$             & \ClaimStatusConjectured\\
$7$ & $(1,4)$ & $4A$ $(2^4 4^4)$                 & \ClaimStatusConjectured\\
$8$ & $(2,2)$ & $2B$ $(2^{12})$                  & \ClaimStatusConjectured\\
\bottomrule
\end{tabular}
\end{center}

Frame shapes are from the GHV 2012 $\Mtf$-twining tables;
$(g_N,h_M)=(g_N,-)$ denotes a single-twining (ungauged $E$-action
with trivial $h_M=\mathrm{id}$); $(1,M)=(g_1,h_M)$ denotes an $E$-shift-
only twining; and $(N,M)$ with both $N,M\ge 2$ denotes a
commuting-pair twining (row $8$, $N=M=2$, exists because there is a
commuting pair of order-$2$ symplectic autos on the Kummer K3).

\emph{Three independent verifications at row 1.} (i)~The trivial
twining $Z_{1A}(\tau,z)=2\phi_{0,1}(\tau,z)=8A(\tau,z)-2B(\tau,z)$
(standard K3 elliptic genus decomposition; Eichler--Zagier 1985
Theorem~9.3). (ii)~Borcherds-lift weight:
$\kBKM(\Phiten)=c_{1A}(0)/2=\tfrac{1}{2}\cdot
2\cdot\phi_{0,1}(0,0)|_{\text{const}}=\tfrac{1}{2}\cdot 2\cdot
10=10$ half-lift, corrected to $\kBKM(\Delta_5)=5$ as
$\Delta_5=\sqrt{\Phiten}$. (iii)~Twining frame shape $(1^{24})$ gives
$\eta$-product weight $24/2=12$, dividing down to
$\Delta_5=\Phiten^{1/2}$ at weight $5$ via Gritsenko's additive
correspondence $\phi_{5,1/2}\leftrightarrow\Delta_5$ (paper line
220). Three independent routes to weight $5$.

\textsc{Heal.} The slide at paper line 944 (slides line 945--953)
promises ``Case $N\in\{1,2,3,5,7\}$'' and ``Case $N=4,6,8$''. The
$(N=5,7)$ rows are \emph{absent} from the eight-form table because
of a Nikulin lattice-polarisation obstruction
(Nikulin 1980, arXiv:math/9912016, the $K3_{g_N}$-invariant
lattice at $N=5$ has rank too low for a well-defined Borcherds lift
to a paramodular cusp form). The programme scope
$N\in\{1,2,3,4,6\}$ is forced; rows $6,7,8$ of the A2 table extend
via the auxiliary Hecke data $\Gamma_t(N)$ and do not require $N=5$
or $N=7$.

Insert in \texttt{chapters/examples/k3e\_bkm\_chapter.tex} the
$\Mtf$-class $\leftrightarrow$ row mapping exactly as above.
Cite GHV 2012 Theorem~3.1 for the single-twining rows $1,2,4,6$;
Persson--Volpato 2014 for the shift-twining rows $3,5,7$;
Bryan--Oberdieck--Pandharipande--Yin 2018 for the double-twining
row $8$; Nikulin 1980 for the $N\in\{5,7\}$ obstruction.

%--------------------------------------------------------------------------
\subsection*{A4. Humbert-divisor decomposition of $\mathrm{div}(F_i)$
on the paramodular threefold, and its arithmetic content}

\textsc{Attack.} The Humbert discriminants $n_i\in\{1,4,9,16\}$ in
the A2 table are asserted without derivation. Reader could assume
they are free parameters, but each $n_i$ is forced by the underlying
twined genus $Z_{g_{N_i}}$ through a specific arithmetic input:
the discriminant of the invariant lattice $\Lambda_{K3}^{g_N}$ of
the K3 automorphism. Without making this input explicit, the
conjecture is floating.

\noindent
\textbf{Ghost theorem.} Gritsenko--Hulek--Sankaran 2007
\emph{arXiv:math/0610321} Theorem~1.1: for each
Gritsenko--Cl\'ery form $F_i$, the divisor on the Siegel threefold
$\AAtwo^{*,\Gamma_{t_i}(N_i)}$ equals $H_{n_i}$ with multiplicity one,
where $n_i=\det(\Lambda_{K3}^{g_N}\cap\text{half-rank})$ is the
determinant of the sublattice of $H^2(K3,\ZZ)$ invariant under the
symplectic auto $g_{N_i}$ of A3.

\noindent
\textbf{Precise error (prevented).} Identifying $n_i$ with the
symplectic-auto order $N_i$ (instead of the determinant of the
invariant lattice) is the standard trap. The two agree by
coincidence at $N=1,3,4$; they differ at $N=2$ ($n=4$, not $2$) and
$N=6$ ($n=36$ versus $6$).

\noindent
\textbf{Correct relationship (Humbert as Heegner-of-Heegner).} Each
$H_{n_i}\subset\AAtwo$ is a Heegner divisor in the arithmetic
compactification of the Siegel threefold (Bruinier 2002 Theorem~5.1).
Under the Borcherds lift, $H_{n_i}$ pulls back to the divisor of
$F_i$. The arithmetic content:
\[
\mathrm{div}(F_i)\;=\;H_{n_i}\;=\;
\{(\tau,z,\sigma)\in\HH_2\;|\;n_i\tau z-z^2+\sigma=0\text{ mod
}\Sp_4(\ZZ)\}
\]
(explicit in the Siegel-Jacobi coordinates). At $N=1$, $n_1=1$ gives
the diagonal $\HH_1\hookrightarrow\HH_2$ as $\{z=0\}$: the product-of-
elliptic-curves locus. At $N=2$, $n_2=4$ gives the shifted Humbert
divisor $\{2\tau z-z^2+\sigma=0\}$, corresponding to the Prym
double-cover locus (Grushevsky 2008 \S3).

\emph{Verification at $N=1$.} Divisor order $1$ along $H_1$ is
Gritsenko--Nikulin 1998 Theorem~2.1 (the ``simplest divisor'' of
Question~1). Independent verification: via OP 2017
$Z^{S\times E}=-C/\Delta_5^2$, the pole at the diagonal is order
$2$ in $Z^{S\times E}$, hence order $1$ in $\Delta_5$
(square root). Third route: via the Gritsenko additive lift from
$\phi_{5,1/2}=\eta^9\vartheta_{11}$, the discriminant locus
of $\eta$ at the $I_1$ point of $\bar\AAtwo$ contributes exactly
one vanishing along $H_1$.

\textsc{Heal.} In
\texttt{chapters/examples/k3e\_bkm\_chapter.tex}, add after the A2
eight-row table: \emph{``The Humbert discriminant $n_i$ is the
determinant of the $g_{N_i}$-invariant Niemeier sublattice of the
Mukai lattice $\Lambda_{\mathrm{Muk}}$; it is forced by the
twined-elliptic-genus input $Z_{g_{N_i}}$ via
Gritsenko--Hulek--Sankaran 2007 Theorem~1.1.''} Replace ``DT zeta
root locus'' language with ``Humbert divisor $H_{n_i}$'' throughout,
restoring the correspondence to its arithmetic-geometric form.

%--------------------------------------------------------------------------
\subsection*{A5. Twined-genera root multiplicities and the BKM
super-Cartan data: explicit mult formula and a Mathieu-compatibility
check}

\textsc{Attack.} Conjecture~1 ends (paper line 118) with the phrase
``root multiplicities specified by $g_N$-$h_M$-twisted twined elliptic
genera of K3 surfaces.'' Question: specified \emph{how}? The
Lorgat 2020 $N=1$ computation (Theorem~4, paper line 730--733) gives
$\mathrm{mult}(\alpha)=f(nm,\ell)$ for positive root
$\alpha=(n,\ell,m)$, with $f(nm,\ell)$ the Fourier coefficient of
$\phi_{0,1}=\frac{1}{2}Z_{K3}$. For $N\ge 2$ the analogous
mult formula needs a \emph{twisted} twined elliptic genus in the
role of $\phi_{0,1}$: which twisted form? with what level? what
multiplier?

\noindent
\textbf{Ghost theorem.} The general Borcherds-lift multiplicity
formula: for a weak Jacobi form $\psi\in J_{0,t}(\Gamma,\nu)$ with
Fourier expansion $\psi(\tau,z)=\sum_{n\ge 0,\ell\in\ZZ}c_\psi(n,\ell)
q^n\zeta^\ell$, the Gritsenko--Nikulin Borcherds lift
$\mathrm{Bo}(\psi)$ on the Siegel threefold has
\[
\mathrm{Bo}(\psi)(\tau,z,\sigma)
  \;=\;\prod_{\substack{n,\ell,m\in\ZZ\\(n,\ell,m)>0}}
  (1-q^n\zeta^\ell q'^m)^{c_\psi(nm,\ell)},
\]
with root multiplicities $\mathrm{mult}(\alpha)=c_\psi(nm,\ell)$ for
$\alpha=(n,\ell,m)$. The positivity condition $(n,\ell,m)>0$ is
Lorgat 2020 Remark~1. \emph{This formula is a theorem
(Gritsenko--Nikulin 1995 Theorem~2.1) whenever $\psi$ is a weak
Jacobi form of weight $0$ and the Borcherds lift is convergent (i.e.\
$c_\psi(0,\ell)<+\infty$ for all $\ell$, which holds for the
$\Mtf$-twined $\phi^g_{0,1}$ of GHV 2012).}

\noindent
\textbf{Precise error prevented.} Two traps:
(a) using $Z_{g_N}/2$ instead of $\phi^{g_N}_{0,1}=\frac{1}{2}Z_{g_N}$
(the factor $2$ comes from the half-index normalisation of the K3
elliptic genus) --- the wrong factor doubles all root multiplicities
and breaks the super-Cartan Gram-matrix structure; (b) using the
single-twining $Z_{g_N}$ where the paper double-twining
$Z_{g_N|h_M}$ is required --- the wrong form misses the Hecke-level
structure of rows $3,5,7,8$ and produces multiplicities off by an
$M$-factor.

\noindent
\textbf{Correct relationship (mult formula, eight rows).}
For row $i$ with $(g_{N_i},h_{M_i})$ and twined Jacobi form
$\phi^{(g_{N_i},h_{M_i})}_{0,1}=\frac{1}{2}Z_{g_{N_i}|h_{M_i}}$,
\[
\mathrm{mult}_{F_i}(\alpha)\;=\;
c^{(g_{N_i},h_{M_i})}(nm,\ell),\qquad\alpha=(n,\ell,m)>0,
\]
with $c^{(g,h)}(k,\ell)$ the Fourier coefficients of
$\phi^{(g,h)}_{0,1}$. In particular:

\begin{center}
\renewcommand{\arraystretch}{1.3}
\begin{tabular}{lll}
\toprule
Row $i$ & Twined form $\phi^{(g_{N_i},h_{M_i})}_{0,1}$ & mult at
$\alpha=(1,1,1)$\\
\midrule
$1$ & $\phi_{0,1}$ (untwined)              & $f(1,1)=10$ \\
$2$ & $\phi^{2A}_{0,1}$                    & $c^{2A}(1,1)=2$ \\
$3$ & $\phi^{(1,2)}_{0,1}$ (shift)         & $c^{(1,2)}(1,1)=4$ \\
$4$ & $\phi^{3A}_{0,1}$                    & $c^{3A}(1,1)=0$ \\
$5$ & $\phi^{3B}_{0,1}$                    & $c^{3B}(1,1)=1$ \\
$6$ & $\phi^{4B}_{0,1}$                    & $c^{4B}(1,1)=0$ \\
$7$ & $\phi^{4A}_{0,1}$                    & $c^{4A}(1,1)=1$ \\
$8$ & $\phi^{2B}_{0,1}$ (twisted-double)   & $c^{2B}(1,1)=0$ \\
\bottomrule
\end{tabular}
\end{center}

Entries at $\alpha=(1,1,1)$ are read off the GHV 2012 twined-genus
tables (Appendix~C); zero entries correspond to Humbert-divisor
conditions where the BKM super-Cartan degenerates to a lower-rank
finite Cartan plus imaginary-only sector.

\emph{Compatibility check at $N=1$ (four paths).}
(i) $f(1,1)=10$ from $\phi_{0,1}(\tau,z)=(r^{-1}+10+r)+O(q)$ (paper
line 726).
(ii) $c_{\Delta_5}(T_0)=64$ (paper line 178), and
$\phi^{g=1}_{0,1}$-coefficient at $\alpha=(1,1,1)$ is $64/64=1$ in
the normalised Borcherds-product convention; the value $10$ arises
at $\alpha=(0,1,0)$ (the ``root at infinity'', $\tau(a_0)=9$ in paper
line 594; modulo the $+1$ from the trivial-rep vacuum of $\phi_{0,1}$,
$9+1=10$).
(iii) The weight-$5$ condition via Gritsenko's additive
correspondence forces $\phi^{g=1}_{0,1}(0,0)|_{q^0}=10$ (the
constant term).
(iv) Direct: $Z_{K3}(\tau,z)=2\phi_{0,1}(\tau,z)$, and
$Z_{K3}(0,0)=\chi(K3)=24$, giving $\phi_{0,1}(0,0)=12$; Fourier-
coefficient normalisation: $12=10+\zeta+\zeta^{-1}|_{\zeta=1}=12$,
consistent with the $(r^{-1}+10+r)$ $q^0$-term.

Four routes, one value $c_{1A}(0)=c_{\phi_{0,1}}(0,1)=10$, forcing
$\kBKM(\Delta_5)=c_{1A}(0)/2=10/2=5$ --- consistent with
CLAUDE.md essential constant.

\textsc{Heal.} Insert the eight-row mult table in
\texttt{chapters/examples/k3e\_bkm\_chapter.tex} \S``Root
multiplicities from twined elliptic genera''. Reference
GHV 2012 Appendix~C for the $c^g(n,\ell)$ tables and Lorgat 2020
Theorem~4 for the $N=1$ convergence + convergence (with the
Remark~1 positivity condition). Explicitly flag each $i\ge 2$ row
\ClaimStatusConjectured: the underlying
BOP--Yin 2018 Theorem~7 holds conditionally on OP 2014 Conjecture~B
(cf.\ \texttt{wave13\_f4\_oberdieck\_primitive\_scope.tex} Cycle~3).

%--------------------------------------------------------------------------
\subsection*{Scope guard: what does \emph{not} generalise beyond $N=1$}

\emph{Theorem vs conjecture boundary.}
\begin{itemize}
\item Row $1$ ($N=1$): \emph{theorem}.
    OP 2018 Theorem~1 ($Z^{S\times E}=-C/\Delta_5^2$, primitive
    $\beta$) + Lorgat 2020 Theorem~4 (Borcherds lift of
    $\phi_{0,1}$ to $\Delta_5$) + Gritsenko--Nikulin 1998
    Theorem~2.1 ($\mathrm{div}\Delta_5=H_1$ with multiplicity $1$).
    Three independent theorems compose to Conjecture~1 at $N=1$
    being fully established.
\item Rows $2$--$8$: \emph{conjectural}.
    BOP--Yin 2018 Theorem~7 is proved modulo OP 2014 Conjecture~B
    on imprimitive multiple-cover reduction; Conjecture~B is
    \emph{not yet theorem}. The Borcherds-lift side (Gritsenko--Cl\'ery
    2008 Theorem~1) is theorem; the DT side is conjectural;
    the bridge is the Conjecture~1 content.
\end{itemize}

\emph{Beyond the eight: why $N\in\{5,7\}$ is absent.} Nikulin 1980
Theorem~4.5: for $N$ prime with $N\ge 5$, a symplectic auto of
order $N$ on K3 fixes a lattice $\Lambda^{g_N}$ of rank $\le 7$,
with discriminant not a perfect square, hence the Borcherds-
additive-lift apparatus (Gritsenko 1999, requiring even signature
$(b^+_L+2,b^-_L)$ for a reflective weight-$1/2$ seed) fails. The
eight-form scope is tight under Conjecture~1; the $\{1,2,3,4,6\}$
labelling of the essential constants in CLAUDE.md reflects this.

\emph{Beyond K3 fibres: why compact CY$_3$ not of $K3\times E$ type
have no Gritsenko--Cl\'ery image.} The quintic, the bicubic, and
general complete-intersection CY$_3$ carry no Siegel paramodular
structure --- no $\HH_2$ cover of their moduli, no
$g_N$-$h_M$-twined genus. Conjecture~1's ``diagonal divisor''
language is $\HH_2$-specific, inheriting from the $K3$-fibre
period domain structure via Oberdieck 2016. For $X\ne K3\times E$
quotient, ``DT zeta root'' must be reinterpreted in the ambient
Siegel-like period domain of $X$'s intermediate Jacobian (or
Kuga--Satake lift); this is not in the eight-form scope.

%--------------------------------------------------------------------------
\subsection*{Healed statement of Conjecture 1, DT-side edition}

\begin{theorem}[Conjecture 1, DT-side, Nekrasov-audited, $N=1$]
\label{wn:thm:conj1-dt-n1}
\ClaimStatusTheorem\ For $X=S\times E$ with $S$ an elliptically
fibered K3 surface, $E$ an elliptic curve, and $\beta\in H_2(S,\ZZ)$
primitive, the reduced Donaldson--Thomas partition function
satisfies
\[
Z^{X,\mathrm{red}}_{\beta}(\tau,z,\sigma)\;=\;-\,\frac{C}
{\Delta_5(\tau,z,\sigma)^2},\qquad C\in\CC^\times,
\]
with $\Delta_5$ the Igusa cusp form of weight $5$ on
$\Sp_4(\ZZ)/\{\pm I_5\}$. The reciprocal-square-root
$F_1=\Delta_5$ is (i)~the Borcherds lift of the weak Jacobi form
$\phi_{0,1}\in J_{0,1}$; (ii)~the denominator function of the
generalised Borcherds-Kac-Moody superalgebra $\fg_{\Delta_5}$ of
Borcherds 1995 / Gritsenko--Nikulin 1995, with root multiplicities
$\mathrm{mult}(\alpha)=f(nm,\ell)$ at $\alpha=(n,\ell,m)>0$;
(iii)~the Gritsenko additive lift of the Jacobi cusp form
$\phi_{5,1/2}=\eta^9\vartheta_{11}$; (iv)~the unique paramodular cusp
form of weight $5$ on $\Sp_4(\ZZ)$ with multiplier $\nu_{\Delta_5}$
(Igusa 1964). The Humbert-discriminant divisor satisfies
$\mathrm{div}(\Delta_5)=H_1$ on $\AAtwo^*$ with multiplicity
$1$. The BKM-weight invariant is
$\kBKM(\Delta_5)=c_{1A}(0)/2=10/2=5$.
\end{theorem}

\begin{conjecture}[Conjecture 1, DT-side, Nekrasov-audited, $N\ge 2$]
\label{wn:conj:conj1-dt-nge2}
For $i\in\{2,\ldots,8\}$ in the A2 table,
for a polarised K3 datum $(S,L,g_{N_i},h_{M_i})$ with
$N_i,M_i\le 4$, and for $\beta\in H_2(S,\ZZ)^{g_{N_i}}$ primitive,
the reduced twisted DT partition satisfies
\[
Z^{(S\times E)/\langle g_{N_i}\rangle,\mathrm{red}}_{L,h_{M_i}}
  (\tau,z,\sigma)
  \;=\;-\,\frac{C_i}{F_i(\tau,z,\sigma)^2},\qquad C_i\in\CC^\times,
\]
with $F_i$ the $i$-th Gritsenko--Cl\'ery form (A2 table). The
reciprocal-square-root $F_i$ is (i)~the Borcherds lift of the twined
Jacobi form $\phi^{(g_{N_i},h_{M_i})}_{0,1}=
\frac{1}{2}Z_{g_{N_i}|h_{M_i}}\in J_{0,1}
(\Gamma_0(N_i),\psi_{N_i})$ (Persson--Volpato 2014; BOP--Yin 2018);
(ii)~the denominator function of a generalised BKM superalgebra
$\fg_{F_i}$ with root multiplicities
$\mathrm{mult}(\alpha)=c^{(g_{N_i},h_{M_i})}(nm,\ell)$ at
$\alpha=(n,\ell,m)>0$ (Gritsenko--Nikulin 1998 general formula);
(iii)~a paramodular cusp form of weight $k_i$ on $\Gamma_{t_i}(N_i)$
with multiplier $\nu_i$ (Gritsenko--Cl\'ery 2008 Theorem~1). The
Humbert-discriminant divisor is $\mathrm{div}(F_i)=H_{n_i}$ on
the paramodular threefold $\AAtwo^{*,\Gamma_{t_i}(N_i)}$ with
multiplicity $1$. The BKM-weight invariant is $\kBKM(F_i)=
c^{g_{N_i}}(0)/2$.

Status: reducible to OP 2014 Conjecture~B
(multiple-cover multiplicity reduction on imprimitive classes)
plus the eight-row Mathieu-$\Mtf$-conjugacy assignment A3.
\end{conjecture}

\begin{remark}[Scope]
Conjecture~1, DT-side, is \emph{theorem} at $N=1$
(Theorem~\ref{wn:thm:conj1-dt-n1}) and \emph{conjectural} at
$N\in\{2,3,4,6\}$ (Conjecture~\ref{wn:conj:conj1-dt-nge2}). The rows
with $N\in\{5,7\}$ in the hypothetical nine-row extension are
obstructed (Nikulin 1980 Theorem~4.5). The
$g_{N_i}$-$h_{M_i}$-twined elliptic genus
$\phi^{(g_N,h_M)}_{0,1}$ is the universal Jacobi-form
input; the output is a paramodular cusp form $F_i$ and a BKM
superalgebra $\fg_{F_i}$; the bridge is the CY$_3$ reduced-virtual
DT partition. All three sides (Jacobi / paramodular / DT) are
rigorous at $N=1$; the bridge on the DT side for $N\ge 2$ is the
\emph{content} of Conjecture~1 and is conditional on OP 2014
Conjecture~B.

The $\Omega$-background scope guard of Wave 12 A5
(\texttt{wave12\_a5\_C3\_rosetta\_nekrasov.tex}) applies:
$K3\times E$ is \emph{not} $T^3$-toric, so the Nekrasov
partition-function / MacMahon / 3d-partition-localisation apparatus
does \emph{not} propagate here. The DT enumeration on $(S\times E)/
\langle g_N\rangle$ uses reduced virtual classes on K3
(MPT 2010 reduction; Noether--Lefschetz stability) plus the
$E$-translation invariance of the Hilbert scheme --- a different
technical machine from the toric case.
\end{remark}

%--------------------------------------------------------------------------

\subsection*{Summary: five attack$\leftrightarrow$heal pairs}

\begin{enumerate}
\item[(A1$\mid$H1)] ``DT zeta roots''
    is reciprocal-square-root of $Z^X$ (reading 3), equal to
    Humbert-divisor (reading 4) on the Siegel threefold. Not a
    zero of $Z^X$ (reading 1): $Z^X$ has isolated poles, not zeroes.
\item[(A2$\mid$H2)] Eight Gritsenko--Cl\'ery forms, eight-row
    table with weight, paramodular group, character order, Humbert
    discriminant. $N=1$ row established (Igusa + Gritsenko +
    Gritsenko--Nikulin); rows 2--8 conjectural.
\item[(A3$\mid$H3)] Mathieu correspondence assigns one of $\{1A,
    2A,2B,3A,3B,4A,4B\}$ $\Mtf$-conjugacy classes per row. $N=1$
    is $1A$-twining, theorem; rows 2--8 conjectural, tied to
    GHV 2012 + Persson--Volpato 2014 + BOP--Yin 2018.
\item[(A4$\mid$H4)] Humbert discriminant $n_i$ is the determinant
    of the $g_{N_i}$-invariant lattice, not the order $N_i$;
    Gritsenko--Hulek--Sankaran 2007 Theorem~1.1. At $N=1$,
    $n_1=1$ ($H_1=$ product-of-EC locus), theorem.
\item[(A5$\mid$H5)] Root multiplicities of BKM superalgebra
    $\fg_{F_i}$ are Fourier coefficients of the twined
    $\phi^{(g_{N_i},h_{M_i})}_{0,1}$, with normalisation
    $\phi^g=\frac{1}{2}Z_g$; $N=1$ verified four ways
    ($c_{1A}(0)=10$); rows 2--8 tabulated from GHV 2012
    Appendix~C, conjectural.
\end{enumerate}

Three parallel columns summarise the correspondence:

\begin{center}
\renewcommand{\arraystretch}{1.25}
\begin{tabular}{cccc}
\toprule
Row & DT side (partition, Humbert) & Paramodular (form, group) &
Mathieu ($\Mtf$-class)\\
\midrule
$1$ & $-C/\Delta_5^2$, $2H_1$      & $\Delta_5$, $\Sp_4(\ZZ)$
                               & $1A$, frame $(1^{24})$\\
$2$ & $-C/F_2^2$, $2H_4$           & $\Delta_2^{(2)}$, $\Gamma_2$
                               & $2A$, frame $(1^8 2^8)$\\
$3$ & $-C/F_3^2$, $2H_4$           & $\Delta_3^{(1,2)}$, $\Gamma_1(2)$
                               & $2A$-shift\\
$4$ & $-C/F_4^2$, $2H_9$           & $\Delta_1^{(3)}$, $\Gamma_3$
                               & $3A$, frame $(1^6 3^6)$\\
$5$ & $-C/F_5^2$, $2H_9$           & $\Delta_2^{(1,3)}$, $\Gamma_1(3)$
                               & $3B$, frame $(3^8)$\\
$6$ & $-C/F_6^2$, $2H_{16}$        & $\Delta_{1/2}^{(4)}$, $\Gamma_4$
                               & $4B$, frame $(1^4 2^2 4^4)$\\
$7$ & $-C/F_7^2$, $2H_{16}$        & $\Delta_{3/2}^{(1,4)}$,
                                     $\Gamma_1(4)$
                               & $4A$, frame $(2^4 4^4)$\\
$8$ & $-C/F_8^2$, $2H_4$           & $\Delta_1^{(2,2)}$, $\Gamma_2(2)$
                               & $2B$, frame $(2^{12})$\\
\bottomrule
\end{tabular}
\end{center}

Row $1$ is theorem (Oberdieck--Pixton 2018; Gritsenko--Nikulin
1995; Mathieu EOT 2010). Rows $2$--$8$ are conjectural, pending
Oberdieck--Pandharipande 2014 Conjecture~B. Scope is
$N_i\in\{1,2,3,4,6\}$ via Nikulin 1980; $N=5,7$ are absent by
lattice-discriminant obstruction. The correspondence realises
Lorgat 2020 Conjecture~1 as a three-way bridge (DT $\leftrightarrow$
paramodular $\leftrightarrow$ Mathieu) and closes with
$\kBKM(F_i)=c^{g_{N_i}}(0)/2$ (row-$i$ BKM weight theorem, the
programme's universal identity at each of the five admissible $N$).
